\section{\sys Design}

\subsection{Design Principles}

\sys follows three key design principles to address the limitations of existing GPU stacks:

\paragraph{Principle 1: Decouple stable mechanisms from dynamic programmable policies}
\noindent\textbf{Our Approach:} \sys explicitly adopts a principle of decoupling stable mechanisms from dynamically programmable policies. We strictly isolate the stable, minimal set of resource management primitives from implementation complexity into clearly defined and stable mechanism APIs exposed by the \sysdriver and runtime for device code. Policy logic is expressed separately as dynamically loadable eBPF programs, attaching to these stable mechanism hooks. This design trades off initial complexity in defining and implementing clear mechanism interfaces but gains substantial long-term flexibility and adaptability: new policies can be rapidly developed, experimented with, and deployed without invasive changes to underlying drivers or runtime components, greatly enhancing GPU datacenter agility and maintainability.

\paragraph{Principle 2: Unified CPU--GPU policy programmability interface and control plane}
\sys explicitly introduces a unified and symmetric CPU--GPU policy programmability paradigm. Rather than relying on disparate, incompatible policy frameworks on host and device sides, we mandate a single, consistent policy representation: all programmable logic across CPU and GPU boundaries is expressed using the same restricted subset of eBPF, with the same abstraction like virtual memory address and state management interface. Policies undergo unified static verification and execution models on both sides, enabling seamless migration, partitioning, and coordination of policy logic between host-side and device-side execution points. Although this unified paradigm intentionally restricts certain aspects of native GPU-side programmability (e.g., direct CUDA kernel manipulations), it provides substantial practical benefits in eliminating policy fragmentation, enabling straightforward cross-layer coordination, and greatly simplifying the development and deployment of adaptive GPU datacenter resource management policies.

\paragraph{Principle 3: Safety and efficiency through verified bounded execution and interfaces}
\noindent\textbf{Our Approach:} \sys explicitly prioritizes both safety and efficiency through a principle of verified bounded execution and interfaces. Rather than allowing arbitrary programmable logic, we restrict all policies—whether executed within GPU kernels or \sysdriver hooks—to a rigorously limited eBPF instructions and interfaces. A unified static verifier ensures that these policies are memory-safe, free from unbounded loops, guaranteed to terminate deterministically, and strictly bounded in resource usage and behavior. A lightweight runtime further enforces stringent per-invocation compute and memory limits, and employs GPU-specific optimizations to minimize overhead when executing device-side instrumentation. This design explicitly trades off some expressiveness in policy programming to ensure the reliability, determinism, and minimal latency overhead required in modern production GPU applications.

\subsection{High-Level Architecture}
\label{sec:arch}

\begin{figure}
\centering
\fbox{\parbox{0.8\columnwidth}{\centering [Figure: \sys Design]}}
\vspace{-.5ex}
\caption{\sys high-level design: unified eBPF runtime across CPU and GPU.}
\label{fig:system-design}
\end{figure}

Figure~\ref{fig:system-design} shows the high-level architecture of \sys, including the CPU--GPU boundary, \sysdriver extensions, and GPU kernel-level runtime integration points. At its core, \sys provides a unified eBPF runtime that manages verification and just-in-time (JIT) compilation in the operating system kernel or the userspace loader.

Policy authors write \sys programs using standard eBPF tooling (e.g., \texttt{clang} with \texttt{libbpf}, \texttt{bpftool}, or \texttt{bpftrace}). Programs are compiled to eBPF bytecode and may access eBPF maps that \sys places in memory regions visible to both the OS kernel and GPU kernels. The bytecode is loaded into the kernel, where it undergoes standard eBPF verification to ensure memory safety, bounded execution, and adherence to the helper API.

Once verified, the same bytecode can be installed in two ways:

\begin{itemize}[leftmargin=*]
  \item \emph{Host side:} \sys attaches eBPF programs to hooks in \sysdriver at critical policy decision points in the GPU stack (e.g., memory prefetching, migration, and kernel scheduling). These hooks follow a \texttt{struct\_ops}-style callback interface: when the driver reaches a hook, it invokes the attached eBPF handler with an event-specific context structure. The handler returns a compact decision (e.g., an action code or priority) that the driver enforces or using safe kfunc calls into the driver code.
  \item \emph{Device side:} for execution inside GPU kernels, \sys translates the verified eBPF bytecode to device bytecode and injects it into the target GPU context. Lightweight trampolines inserted at kernel entry or other selected sites save minimal live state, call into the translated eBPF snippet, and then restore state to resume the original control flow. Device-side eBPF programs update the same maps as their host-side counterparts.
\end{itemize}

A GPU-side sandbox runtime enforces strong safety guarantees: it bounds the number of instructions per invocation, carefully manages register usage to avoid corrupting application state, and restricts memory accesses to designated shared regions (e.g., eBPF maps and scratch space). Each injected program therefore operates in a protected domain that cannot interfere with application kernels, while still enabling adaptive policy decisions on critical GPU execution paths.

Shared eBPF maps provide efficient state coordination between \sysdriver and eBPF policy logic on the host, and between host and GPU policies across the CPU--GPU boundary. This design yields real-time, low-overhead cross-layer observability and control while retaining the safety guarantees of the eBPF framework.

\subsection{\sys Programming Model}
\label{sec:prog-model}

A \sys policy is a distributed application consisting of eBPF code that runs in
privileged contexts on the CPU and GPU together with a user-space control-plane
component. All pieces program against a unified abstraction: a restricted eBPF IR,
a small set of attach points, and shared eBPF maps and metadata.

\subsubsection{Policy Code (CPU \& GPU)}

\sys policies are expressed as eBPF programs using a unified intermediate representation: a restricted subset of eBPF instructions and helpers. Each program defines multiple entry points corresponding to attach points on either the host or GPU device. On the host, verified eBPF bytecode executes directly in privileged kernel contexts provided by \sysdriver, allowing policies to respond to critical memory events (such as device memory misses, eviction candidate selection) and kernel scheduling decisions (kernel submission, admission, preemption, completion). On the GPU, the same verified bytecode is translated to PTX instructions and executed through lightweight trampolines injected at kernel entry points, selected memory operations, and synchronization or phase boundaries. In both host and GPU contexts, policy handlers receive structured event contexts containing metadata (e.g., virtual addresses, region IDs, kernel IDs, warp identifiers, and resource usage metrics), and return compact policy decisions enforced directly by underlying GPU mechanisms.

\subsubsection{Shared State \& Metadata}

To enable coherent policy decisions across the CPU-GPU boundary, \sys provides a unified state abstraction via shared eBPF maps, placed in unified memory visible to both CPU and GPU policy code. Maps store policy-relevant state, including counters, histograms, per-region hotness metrics, per-tenant scheduling priorities, and runtime configuration parameters. This shared state allows device-side probes, such as GPU memory access instrumentation or warp progress signals, to update policy metadata that host-side memory or scheduling hooks later consult. \sys transparently propagates and translates core metadata—such as virtual memory addresses, kernel identifiers, tenant IDs, timestamps, and execution contexts—across CPU and GPU, presenting a logically unified identifier space. This design allows policy developers to implement sophisticated cross-layer logic (e.g., GPU-side access tracking coupled with host-side eviction or scheduling) without manually handling address translation or cross-boundary synchronization complexities.

\subsubsection{Control-Plane Application and Lifecycle}

A \sys policy also includes a dedicated user-space control-plane component responsible for lifecycle management, dynamic reconfiguration, and policy introspection. This component loads compiled eBPF bytecode into the OS kernel, instantiates and manages eBPF maps, and attaches policy sections to defined host and GPU device hook points exposed by \sysdriver and GPU-side runtime. It enables dynamic reconfiguration of deployed policies at runtime by updating parameters (e.g., thresholds, sampling rates, eviction aggressiveness, power caps) stored in shared maps, and facilitates atomic swapping or updating of attached policy logic without restarting applications or kernels. Furthermore, the control-plane application continuously monitors policy operation by reading metrics and state from maps—such as memory hotness counters, per-tenant scheduling metrics, and queueing latencies—and exports this information to external observability frameworks or adaptation controllers. Thus, from a developer’s perspective, a complete \sys policy application comprises policy logic for both CPU and GPU execution contexts along with a complementary control-plane component orchestrating policy behavior and evolution over time.

\subsection{Safety Model and Verification}
\label{sec:safety}

\sys enforces a unified safety model for all policy code running in \sysdriver and
inside GPU kernels. Policies must pass static verification before deployment and
remain within strict runtime bounds when executed on the GPU.

\paragraph{Static verification.}
On the host, \sys reuses the standard eBPF verifier to ensure that programs are
memory-safe, free of unbounded loops, and respect limits on stack depth and helper
usage. The verifier only allows calls to approved helper functions and access to
supported map types, and rejects programs that could corrupt driver state or violate
kernel invariants.

For device-side execution, \sys applies the same verifier to the eBPF bytecode
before any translation to PTX. This ensures that the control-flow and memory-safety
properties proven by the verifier hold for both host and device instantiations of
the policy.

\paragraph{GPU sandbox and bounded execution.}
After verification, device-side policy code is translated to PTX and executed in a
sandboxed runtime. The GPU sandbox bounds the number of instructions per invocation,
constrains register allocation to a reserved set to avoid clobbering application
state, and restricts memory accesses to designated shared regions such as eBPF maps
and per-warp scratch buffers. Lightweight PTX trampolines save and restore only the
minimal live register and predicate state around calls into the translated eBPF code.

Together, static verification and the GPU sandbox guarantee that injected policy
logic is bounded, deterministic, and cannot hang the GPU, corrupt kernel or application
state, or impose unbounded latency on performance-critical paths such as memory
fault handling or kernel admission. This makes it safe to execute policy code
directly at the points where GPU resource management decisions are made.


\subsection{Policy Interface for Memory}
\label{sec:design-mem}

Modern GPU memory subsystems migrate and evict data across device and host memory
using fixed, internally implemented heuristics, often driven by page faults or
simple recency-based policies. These heuristics can introduce latency spikes and
fail to adapt to diverse access patterns. \sys exposes a bidirectional interface
for memory policies: host-side hooks at key decision points and device-side probes
that collect fine-grained access information.

\paragraph{Host-side memory hooks.}
On the host, \sys introduces attach points in the GPU memory management path at
several stages:

\begin{itemize}[leftmargin=*]
  \item \emph{Miss hooks}, triggered when the GPU issues a memory access that
        requires the host to allocate, migrate, or otherwise service data for
        a region.
  \item \emph{Prefetch hooks}, triggered at safe points where a policy may suggest
        regions to migrate proactively based on observed access patterns and
        available capacity.
  \item \emph{Eviction hooks}, triggered when device memory is under pressure and
        the system must select one or more regions to evict or demote.
\end{itemize}

Each hook passes a context structure describing the event (e.g., the accessed
virtual address or region identifier, access type, region size, current residency,
and local memory pressure metrics) to an attached eBPF handler. The handler can
consult shared maps (for example, per-region hotness and tenant budgets) and return
a compact decision: \emph{migrate}, \emph{retain}, \emph{prefetch}, or a ranking
score among candidates. The underlying memory subsystem remains responsible for
correctness (e.g., ensuring some victim is chosen), but delegates policy choices to
the eBPF program.

\paragraph{Device-side memory instrumentation.}
On the GPU, \sys injects lightweight probes around selected memory operation sites
in target kernels. These probes operate at warp granularity, observing access
metadata such as a logical region identifier, warp and SM IDs, and a kernel or
phase identifier. Probes update per-region or per-object statistics in shared maps,
such as access counts, approximate working-set membership, or simple stride or
locality signatures. To bound overhead, \sys supports sampling (e.g., one in every
$N$ accesses) and per-warp aggregation before updating maps.

\paragraph{Example usage: programmable prefetch and eviction.}
Together, these host and device interfaces let \sys support programmable prefetch
and eviction without changing the underlying memory subsystem. Device-side probes
track working sets and access frequencies in shared maps; host-side miss and
prefetch hooks consult this state to prefetch regions that are predicted to be
accessed soon, and eviction hooks select cold regions proactively. Because these
policies are expressed entirely in eBPF over generic hooks and maps, operators can
deploy multiple algorithms (e.g., threshold-based hotness, recency-biased eviction,
tenant-aware placement) and evolve them without modifying driver code or GPU
kernels.

\subsection{Policy Interface for Scheduling}
\label{sec:design-sched}

GPU schedulers typically provide fixed, coarse-grained policies (e.g., FIFO queues
with limited priority control), which are insufficient in multi-tenant environments
with mixed latency-sensitive and batch workloads or power constraints. \sys exposes
scheduling interfaces on both host and device so that policies can implement
fine-grained priority, deadline-aware, or power-aware scheduling.

\paragraph{Host-side scheduling hooks.}
On the host, \sys extends the GPU kernel lifecycle with attach points at:

\begin{itemize}[leftmargin=*]
  \item \emph{Submit}, when a kernel is enqueued by a framework or runtime. The
        context includes kernel metadata (e.g., grid and block size), tenant or
        stream identifier, and any priority hints.
  \item \emph{Admit}, when the scheduler is about to dispatch a ready kernel to
        a GPU context or hardware queue. The policy can re-order ready kernels,
        defer admission, or adjust effective priorities.
  \item \emph{Preempt}, when a preemption opportunity arises, for example at a
        cooperative yield point or at a scheduling quantum boundary. The policy
        can decide whether to continue or preempt a kernel based on its progress
        and the presence of higher-priority work.
  \item \emph{Complete}, when a kernel finishes. Policies use this hook to update
        per-tenant statistics, virtual time, or longer-term admission control state.
\end{itemize}

Attached eBPF schedulers read these contexts and shared maps (e.g., per-tenant
latency and throughput statistics, outstanding work across queues) and return
compact scheduling decisions such as which kernel to run next, whether to throttle
a tenant, or whether to allow preemption.

\paragraph{Device-side scheduling signals.}
On the GPU, \sys injects probes at kernel entry and at selected synchronization or
phase boundaries within kernels. These probes can record per-warp or per-block
progress (e.g., iterations completed in a loop, entry into a barrier-heavy phase),
mark when a kernel enters a phase that is safe to preempt, or export simple signals
such as ``memory-bound'' or ``compute-bound'' to shared maps. Device-side probes
are restricted to bounded, lightweight operations and cannot directly manipulate
hardware queues; instead, they communicate hints and progress information that host
policies consume at the next scheduling hook.

\paragraph{Example usage: fine-grained priority and power-aware scheduling.}
Using these interfaces, \sys supports a variety of scheduling policies. A priority
scheduler maintains per-tenant priorities and latency statistics in maps keyed by
tenant ID. Submit and admit hooks tag and order kernels based on these priorities,
while preempt hooks consult device-side progress signals to preempt long-running
low-priority kernels at synchronization points when high-priority work arrives.
Similarly, a power-aware scheduler combines host-visible power and temperature
readings with device-side activity counters to throttle or reschedule kernels as
the GPU approaches configured power or thermal limits. These policies are expressed
entirely in eBPF and built over generic scheduling hooks, allowing operators to
deploy and evolve them without modifying GPU applications or vendor schedulers.
